\chapter{Практическая часть}
\label{ch:chap2}

\section{Знакомство с ImGUI и ImPlot}

\subsection{ImGui}

ImGui (Immediate Mode Graphical User Interface) — это библиотека и паттерн проектирования для создания быстрых графических 
интерфейсов (UI), где элементы отрисовываются каждый кадр непосредственно в коде. Она не хранит состояние UI, что упрощает 
синхронизацию и позволяет создавать инструменты для отладки, редакторы и интерфейсы для C++ приложений. \\

Приложение создает сцену из графических примитивов, таких как фигуры, линии и т.д., образуя сцены. 
Графическая библиотека хранит модель сцены в памяти. При отрисовки каждого кадра, приложение обращается к графическому 
API за обновлением сцены. Затем библиотека отрисовывает сцену.

\begin{figure}[H]
    \centering
    \includegraphics[width=1.0\textwidth]{imgui.png}
    \caption{Принцип работы ImGui}
\end{figure}

Особенности:

\begin{enumerate}
    \item Интерфейс ImGui вызывается в цикле, каждая итерация цикла считается одним кадром.

    \item Перед началом кадра библиотека должна получить все данные касательно положения мыши, нажатых клавиш и прочих вводных.

    \item Затем вызывается функция, в которой пользователь библиотеки получает структуру ui, рисует свой интерфейс. В этот момент библиотека в реальном времени формирует все виджеты, для каждого виджета определяет попала ли на него мышь, и формирует вершины и треугольники, которые потом будут рисоваться.

    \item Теперь библиотека возвращает все сформированные mesh’и и отдаёт их на рисование бэкенду. Он рисует это на экран.

    \item И так повторяется 60 кадров в секунду (или больше).
\end{enumerate}

\subsection{ImPlot}

ImPlot — это библиотека для быстрого построения 2D-графиков в реальном времени, разработанная специально для Dear ImGui. 
Она использует концепцию непосредственного выполнения (immediate mode), позволяя визуализировать данные с высокой 
производительностью, минимальными затратами ресурсов и легкой интеграцией без внешних зависимостей.

\begin{figure}[H]
    \centering
    \includegraphics[width=1.0\textwidth]{implot.png}
    \caption{Возможности ImPlot}
\end{figure}

\subsection{Работа с ImGui (схематично)}

\begin{figure}[H]
    \centering
    \includegraphics[width=1.0\textwidth]{prog.png}
    \caption{Схема программы с ImGui}
\end{figure}

Пример инициализации:

\begin{lstlisting}
    // 1) Инициализация SDL
    SDL_Init(SDL_INIT_VIDEO | SDL_INIT_TIMER);
    SDL_Window* window = SDL_CreateWindow(
        "Backend start", SDL_WINDOWPOS_CENTERED, SDL_WINDOWPOS_CENTERED,
        1024, 768, SDL_WINDOW_OPENGL | SDL_WINDOW_RESIZABLE);
    SDL_GLContext gl_context = SDL_GL_CreateContext(window);

    // 2) Инициализация контекста Dear Imgui
    ImGui::CreateContext();
    ImPlot::CreateContext();

    // Ввод\вывод
    ImGuiIO& io = ImGui::GetIO(); (void)io;
    io.ConfigFlags |= ImGuiConfigFlags_DockingEnable;         // Включить Docking

    // 2.1) Привязка Imgui к SDL2 и OpenGl backend'ам
    ImGui_ImplSDL2_InitForOpenGL(window, gl_context);
    ImGui_ImplOpenGL3_Init("#version 330");
\end{lstlisting}

\subsection{Основные элементы ImGui и ImPlot}

\subsubsection{Input}

Виджеты, предназначенные для ввода чисел (int, float, double)

\begin{figure}[H]
    \centering
    \includegraphics[width=1.0\textwidth]{input.png}
    \caption{Пример input виджетов}
\end{figure}

\begin{lstlisting}
    static int i0 = 123;
    ImGui::InputInt("input int", &i0);

    static float f0 = 0.001f;
    ImGui::InputFloat("input float", &f0, 0.01f, 1.0f, "%.3f");

    static double d0 = 999999.00000001;
    ImGui::InputDouble("input double", &d0, 0.01f, 1.0f, "%.8f");
\end{lstlisting}

В функции передаются следующие аргументы:

\begin{enumerate}
    \item "input double" - уникальный строковый идентификтор. Отображается возле виджета.
    \item \&d0 - адрес переменной, к которой будет привязан виджет (чье значение будет менять).
    \item 0.01f - минимальное значение.
    \item 1.0f - макисмальное значение.
    \item "\%.8f" - формат вывода.
\end{enumerate}

\subsubsection{Checkbox и radiobutton}

\begin{figure}[H]
    \centering
    
\includegraphics[width=1.0\textwidth]{rb.png}
    \caption{Пример radiobutton и checkbox}
\end{figure}

\begin{lstlisting}
    static bool check = true;
    ImGui::Checkbox("checkbox", &check);

    static int e = 0;
    ImGui::RadioButton("radio a", &e, 0);
    ImGui::RadioButton("radio b", &e, 1);
    ImGui::RadioButton("radio c", &e, 2);
\end{lstlisting}

Radiobutton третьим аргументом принимает значение, которое будет устанавливаться в переменной после нажатия на radiobutton.

\subsubsection{Sliders}

\begin{figure}[H]
    \centering
    \includegraphics[width=1.0\textwidth]{slider.png}
    \caption{Пример слайдеров}
\end{figure}

\begin{lstlisting}
    static int i1 = 0;
    ImGui::SliderInt("slider int", &i1, -1, 3);

    static float f1 = 0.123f, f2 = 0.0f;
    ImGui::SliderFloat("slider float", &f1, 0.0f, 1.0f, "ratio = %.3f");
    ImGui::SliderFloat("slider float (log)", &f2, -10.0f, 10.0f, "%.4f", ImGuiSliderFlags_Logarithmic);
\end{lstlisting}

Синтаксис очень похож на виджеты Input, но есть возможность установить флаги, например, для логарифмической шкалы.

\subsubsection{Plots}

\begin{figure}[H]
    \centering
    \includegraphics[width=1.0\textwidth]{plot.png}
    \caption{Пример графика}
\end{figure}

Код:

\begin{lstlisting}
    if (ImPlot::BeginPlot("I/Q samples", ImVec2(-1, plot_height)))
    {
       ImPlot::SetupAxes("Time", "Amplitude");
       ImPlot::PlotLineG("I component", get_I<int16_t>,
                         &rx_config.rx_samples,
                         rx_config.rx_samples.size());
       ImPlot::PlotLineG("Q component", get_Q<int16_t>,
                         &rx_config.rx_samples,
                         rx_config.rx_samples.size());
       ImPlot::EndPlot();
    }
\end{lstlisting}

Вызвать само окно под график можно с помощью ImPlot::BeginPlot(), куда необходимо передать название графика (это в то же время
его уникальный ID), размер графика в виде структуры ImVec2(width, height). -1 означает, что занимается вся доступная ширина.
Вызов можно удобно делать в if, что предотвращает появление ошибок. \\


Далее идет более детальная настройка графика. С помощью ImPlot::SetupAxes() настраиваются подписи для осей. ImPlot::PlotLineG() позволяет
отобразить данные. Функция прнимает название выводимых данных, сами данные (адрес), callback функцию и размер данных. Callback функция
здесь нужна по той причине, что ImPlot умеет работать только с дефолтными массивами и только с дефолтными типами данных. У меня, например,
была задача вывести вектор комплексных чисел, а это значит, что вектор надо было распарсить на 2 обычных массива с мнимой и реальной
частью. Делать это крайне неудобно, поэтому разработчики ImPlot позаботились об этом заранее и дали возможность передать в функцию
указатель на callback функцию, которая будет автоматически применяться к каждому выводимому объекту. \\

По большому счету пользователю требуется только описать то, как парсить данные, а всю остальную рутинную работу сделает ImPlot.\\

Callback функция пишется в строго определенном стиле:

\begin{lstlisting}
    template <typename T> 
    ImPlotPoint get_I(int idx, void *user_data) {
    auto *vec = static_cast<std::vector<std::complex<T>> *>(user_data);
    return ImPlotPoint(idx, (*vec)[idx].real());
}
\end{lstlisting}

Функция должна обязательно иметь сигнатуру, которая указана выше. Здесь объект ImPlotPoint - точка на графике, idx - индекс, который
ImPlot будет подавать на функцию, user\_data - данные, которые будет подавать ImPlot. Логика такая: ImPlot, допустим, хочет вывести
вектор, он берет нулевой элемент вектора и индекс (0) подставляет вместо idx, а сам элемент вместо user\_data. Дальше идет обработка,
в ходе которой должен быть сформирован и возвращен объект точки, который ImPlot отрисует.


\section{Подводный камень при работе с GUI}
При работе с GUI обязательным является вынос логики работы GUI в отдельный поток, поскольку задачи основного потока могут тормозить
работу GUI и наоборот. Запустить новый поток несложно, но встает вопрос обмена данными между потоками. Для этого использовались 
обычные структуры, которые по ссылки передавались в потоки. Таким образом, например, основной поток принимал с SDR семплы, а 
поток GUI сразу же читал их и выводил, причем в такой схеме не требуется синхронизировать потоки, потому что GUI только читает.

\begin{lstlisting}
     std::thread gui_thread(run_gui, std::ref(tx_config), std::ref(rx_config),
                         std::ref(sdr_config));

  std::thread sdr_thread(SDR_run, std::ref(sdr_config), std::ref(tx_config),
                         std::ref(rx_config));
  gui_thread.join();
  sdr_thread.join();
\end{lstlisting}

\section{Результат}
Имея небольшой функционал, описанный выше, можно построить вполне неплохой GUI:

\begin{figure}[H]
    \centering
    \includegraphics[width=1.0\textwidth]{rx.png}
    \caption{GUI для RX}
\end{figure}

\begin{figure}[H]
    \centering
    \includegraphics[width=1.0\textwidth]{tx.png}
    \caption{GUI для TX}
\end{figure}

\begin{figure}[H]
    \centering
    \includegraphics[width=1.0\textwidth]{sdr.png}
    \caption{GUI для SDR}
\end{figure}

\endinput
